% ── §1.2: Σχετικές Εργασίες ──────────────────────────────────
%       [19]=Liu2021, [29]=McMahan2017, [33]=Yang2019, [36]=Li2020,
%       [37]=Kairouz2021, [41]=Kang2019, [42]=Sarikaya2020,
%       [43]=Yu2020, [44]=Wang2019, [45]=Shapley1953, [46]=Gopinathan2023, [47]=Xu2021

\section{Σχετικές Εργασίες}
\label{sec:related_work}

\subsection{Αναγνώριση Δραστηριότητας Ανθρώπου}

Η αναγνώριση δραστηριότητας μέσω αισθητήρων έχει διανύσει μακρά
ερευνητική πορεία. Οι πρώτες προσεγγίσεις βασίζονταν σε χειροκίνητη
εξαγωγή χαρακτηριστικών (hand-crafted features) και κλασικούς ταξινομητές,
ωστόσο η ευρεία διάδοση των έξυπνων κινητών συσκευών επέτρεψε τη συλλογή
μεγάλων datasets αισθητήρων και τη μετάβαση σε αρχιτεκτονικές βαθιάς
μάθησης. Το καθιερωμένο σύνολο δεδομένων UCI HAR \cite{Anguita2013}
παρέχει δεδομένα επιταχυνσιόμετρου και γυροσκοπίου από 30 εθελοντές για
έξι βασικές δραστηριότητες, θέτοντας ένα ευρέως χρησιμοποιούμενο πρότυπο
αξιολόγησης.

Η εργασία των Hammerla et al.\ \cite{Hammerla2016} έδειξε ότι τα βαθιά
νευρωνικά δίκτυα — ιδίως τα Συνελικτικά Νευρωνικά Δίκτυα (Convolutional
Neural Networks, CNN) και τα δίκτυα Long Short-Term Memory (LSTM) — υπερτερούν
σημαντικά των παραδοσιακών μεθόδων στην ανάλυση σημάτων αισθητήρων. Η
ολοκληρωμένη ανασκόπηση των Chen et al.\ \cite{Chen2021} συγκρίνει
συστηματικά αρχιτεκτονικές CNN, RNN και υβριδικά μοντέλα, αναδεικνύοντας
τα 1D-CNN ως ιδιαίτερα κατάλληλα για HAR λόγω της ικανότητάς τους να
εξάγουν τοπικά χαρακτηριστικά χρονοσειρών αποδοτικά και χωρίς
εκτεταμένη προεπεξεργασία. Παρόμοια ευρήματα αναφέρονται και από τους
Liu et al.\ \cite{Liu2021}, οι οποίοι εφαρμόζουν FL ειδικά για HAR σε
έξυπνα συστήματα υγειονομικής περίθαλψης.

\subsection{Ομοσπονδιακή Μάθηση}

Η Ομοσπονδιακή Μάθηση εισήχθη από τους McMahan et al.\ \cite{McMahan2017}
με τον αλγόριθμο Federated Averaging (FedAvg), ο οποίος επιτρέπει την
αποκεντρωμένη εκπαίδευση μοντέλων με ελάχιστη επικοινωνία μεταξύ server
και clients. Η εννοιολογική θεμελίωση του FL και οι εφαρμογές του
παρουσιάζονται εκτενώς από τους Yang et al.\ \cite{Yang2019}, ενώ οι Li
et al.\ \cite{Li2020} αναλύουν τις βασικές προκλήσεις του πεδίου:
ετερογένεια δεδομένων (Non-IID), αναξιοπιστία clients, και κόστος
επικοινωνίας. Η εκτενής ανασκόπηση των Kairouz et al.\ \cite{Kairouz2021}
αποτελεί σήμερα ένα από τα πλέον αναφερόμενα έργα στο πεδίο, καλύπτοντας
ανοικτά ζητήματα ασφάλειας, ιδιωτικότητας και δικαιοσύνης στο FL.

\subsection{Μηχανισμοί Κινήτρων στο FL}

Η μοντελοποίηση κινήτρων συμμετοχής στο FL μέσω θεωρίας παιγνίων
αποτελεί ενεργό ερευνητικό πεδίο. Οι Kang et al.\ \cite{Kang2019}
πρότειναν έναν μηχανισμό βασισμένο στο παίγνιο Stackelberg για κινητά
δίκτυα, στον οποίο ο server (leader) καθορίζει ανταμοιβές και οι clients
(followers) βελτιστοποιούν την προσπάθειά τους. Οι Sarikaya και Ercetin
\cite{Sarikaya2020} απέδειξαν την ύπαρξη μοναδικής λύσης Nash Equilibrium
στο αντίστοιχο παίγνιο. Οι Yu et al.\ \cite{Yu2020} εισήγαγαν περιορισμούς
δικαιοσύνης (fairness constraints) στον μηχανισμό ανταμοιβής, αποτρέποντας
την εκμετάλλευση των αδύναμων clients από τους ισχυρότερους. Ωστόσο,
καμία από αυτές τις εργασίες δεν συνδυάζει τον μηχανισμό κινήτρων με
συστηματική αξιολόγηση της ποιότητας της συνεισφοράς κάθε client.

\subsection{Αξιολόγηση Συνεισφοράς Clients}

Για την αξιολόγηση της συνεισφοράς, οι Wang et al.\ \cite{Wang2019}
εφάρμοσαν πρώτοι τα Shapley Values \cite{Shapley1953} από τη Συνεργατική
Θεωρία Παιγνίων στο πλαίσιο του FL, αναδεικνύοντάς τα ως θεωρητικά
θεμελιωμένο μέτρο δίκαιης κατανομής ανταμοιβών. Οι Gopinathan et al.\
\cite{Gopinathan2023} επέκτειναν αυτή την προσέγγιση, ενώ οι Xu et al.\
\cite{Xu2021} πρότειναν μια αποδοτική προσέγγιση βασισμένη στη ομοιότητα
gradients (gradient-based Shapley approximation), μειώνοντας δραστικά την
υπολογιστική πολυπλοκότητα. Τα ανωτέρω έργα αντιμετωπίζουν όμως το
πρόβλημα της αξιολόγησης συνεισφοράς ανεξάρτητα από κάποιο μηχανισμό
κινήτρων, χωρίς να κλείνουν τον βρόχο μεταξύ αξιολόγησης και
ανταμοιβής.

\subsection{Κενό στη Βιβλιογραφία}

Από την ανασκόπηση της βιβλιογραφίας προκύπτει ότι, παρά τη σημαντική
πρόοδο σε κάθε έναν από τους ανωτέρω τομείς χωριστά, δεν υπάρχει
ολοκληρωμένη εργασία που να συνδυάζει: (α) FL για HAR σε AAL περιβάλλοντα,
(β) μηχανισμό κινήτρων βασισμένο σε θεωρία παιγνίων, και (γ) δίκαιη
αξιολόγηση συνεισφοράς μέσω Shapley Values. Ο Πίνακας~\ref{tab:related_work}
συγκρίνει τις βασικότερες εργασίες ως προς αυτές τις τέσσερις διαστάσεις.

\begin{table}[H]
\centering
\caption{Σύγκριση σχετικών εργασιών με το προτεινόμενο σύστημα.
(\checkmark = υποστηρίζεται, \textemdash{} = απουσιάζει)}
\label{tab:related_work}
\small
\begin{tabular}{lccccc}
\hline
\textbf{Εργασία} & \textbf{FL} & \textbf{HAR} & \textbf{Κίνητρα} & \textbf{Shapley/LOO} & \textbf{AAL} \\
\hline
Hammerla et al.\ \cite{Hammerla2016}    & ---         & \checkmark  & ---              & ---          & ---          \\
Chen et al.\ \cite{Chen2021}            & ---         & \checkmark  & ---              & ---          & ---          \\
Liu et al.\ \cite{Liu2021}              & \checkmark  & \checkmark  & ---              & ---          & ---          \\
McMahan et al.\ \cite{McMahan2017}      & \checkmark  & ---         & ---              & ---          & ---          \\
Kang et al.\ \cite{Kang2019}            & \checkmark  & ---         & \checkmark       & ---          & ---          \\
Sarikaya \& Ercetin \cite{Sarikaya2020} & \checkmark  & ---         & \checkmark       & ---          & ---          \\
Yu et al.\ \cite{Yu2020}               & \checkmark  & ---         & \checkmark       & ---          & ---          \\
Wang et al.\ \cite{Wang2019}            & \checkmark  & ---         & ---              & \checkmark   & ---          \\
Gopinathan et al.\ \cite{Gopinathan2023}& \checkmark  & ---         & ---              & \checkmark   & ---          \\
Xu et al.\ \cite{Xu2021}               & \checkmark  & ---         & ---              & \checkmark   & ---          \\
\hline
\textbf{Παρούσα εργασία}               & \checkmark  & \checkmark  & \checkmark       & \checkmark   & \checkmark   \\
\hline
\end{tabular}
\end{table}

Η παρούσα εργασία καλύπτει ακριβώς αυτό το κενό — είναι η μόνη που
συνδυάζει και τις τέσσερις διαστάσεις ταυτόχρονα.
